%%%%%%%%%%%%%%%%%%%%%%%%%%%%%%%%%%%%%%%%%%%%%%%%%%%%%%%%%%%%%%%%%%%%%%%%%%%%%%
% 02_formal_setting.tex
%
% Formal setting and assumptions for the analysis of structural diagnosability
% of the enterprise continuum K_EA.
%
% This section fixes:
%   - the OC/ETS structures taken as primitive,
%   - the minimal representational (finiteness) assumption required for
%     decidability statements,
%   - the exact scope of information available to any decision procedure.
%
% No new primitives are introduced.
%%%%%%%%%%%%%%%%%%%%%%%%%%%%%%%%%%%%%%%%%%%%%%%%%%%%%%%%%%%%%%%%%%%%%%%%%%%%%%

We work strictly within the \OC{} core primitives and the enterprise-continuum
specification \KEA{} as fixed by \texttt{ETS.K\_EA.v1.1}. The enterprise continuum
is treated as a formal object equipped with (at minimum) the following immutable
components:
\[
(\Om(\KEA), \dOm(\KEA), \Ax(\KEA), \Pot(\KEA), \Th(\KEA), \Flow(\KEA), \Cyc(\KEA), \kcont(\KEA)).
\]
All symbols above refer exclusively to the corresponding definitions provided
by \OC{} core and \texttt{ETS.K\_EA.v1.1}. No reinterpretation or semantic
extension is permitted.

\subsection{Representational assumption}
To formulate decidability results, we require that an admissible instance of
\KEA{} be provided in a \emph{finite}, machine-readable representation consistent
with \texttt{ETS.K\_EA.v1.1}. This assumption does not introduce a new primitive;
it merely makes explicit the input model required for any effective decision
procedure.

In particular, finiteness applies to the representations of axes, thresholds,
flows, cycles, and all auxiliary data structures prescribed by ETS, yielding a
finite encoding domain for admissible enterprise states.

\subsection{Permitted information for decision procedures}
Any decision procedure considered in this work may inspect only the following
ETS-provided structures of the given \KEA{} instance:
\begin{itemize}
  \item the formal representation of admissible states in \(\Om(\KEA)\) and the
        associated boundary structure \(\dOm(\KEA)\);
  \item the finite sets of axes \(\Ax(\KEA)\) and thresholds \(\Th(\KEA)\);
  \item any explicitly represented flow and cycle data
        \((\Flow(\KEA), \Cyc(\KEA))\) contained in the ETS-conformant input;
  \item the STOP/No-Go constraints and invariants encoded in
        \texttt{ETS.K\_EA.v1.1}.
\end{itemize}
No appeal to managerial intent, KPI observability, or external interpretive
semantics is permitted.

\subsection{Structure of the formal argument}
Within this setting, we proceed as follows: first, structural diagnosability is
defined as a Boolean property determined solely by \OC{} and ETS structures;
second, this property is shown to be well-defined (invariant under
ETS-equivalent representations); and third, decidability is established by
reduction to a finite decision problem over the ETS-conformant representation.

%%% End of file %%%
